\scanpage{116}
\mspage{122}
\deleted{b) Pour tout $p:Y''\unsure{\mathord{-}|\to}Y'$ sur $Y$, si $p$ est une $\unsure{\mathbb D|D}$-équiv. $1$-compl. (resp. $2$-complète), il en est de même de $q:X''\to X'$.}

\deleted{b$_1^{\unsure{*|'}}$) Comme b$_1$), mais en supposant $Y''$ cat. \unsure{ponctuelle}, $Y'$ \deleted{\unsure{d'un}} un objet final}

\deleted{b$_2$) Comme b), mais en supposant $Y''$ cat. ponctuelle, $Y'$ un objet initial.}

\deleted{Je ne \unsure{connais} pas de \unsure{renforcement} co}

La notion de $\mathbb D$-fibration forte a été \deleted{\unsure{prise}} choisie pour \unsure{coiffer} les deux notions duales de \unsure{$\mathbb D$-fib.|$f$-fib.} coh. et hom. (à défaut d'une meilleure notion), tout en restant stable par composition. On pourrait exiger que pour tt

\[Y''\xrightarrow{p}Y'\]

sur $Y$, avec $Y',Y''$ \unsure{Hot-asphériques|tlt-asphériques}, $q:X''-X'$ soit une $\mathbb D$-équivalence $1$-complète (\unsure{voire} même, \deleted{\unsure{une}} $2$-complète), appeler ça les $\mathbb D$-\uline{fibrations} $1$-\uline{complètes} (resp. $2$-\uline{complètes}). \deleted{\unsure{Ça coiffe}} Renforce \deleted{\unsure{quel}} de façon assez élégante la notion de $\mathbb D$-fibration forte, mais a le grand tort de ne pas être stable par composition (que je sache !) ; sauf dans le cas particulier où on a des \unsure{Hom} internes \unsure{sympa} (cf bas de la page 113), auquel cas la notion de $\mathbb D$-équiv. coh. implique déjà celle de $\mathbb D$-équiv. $1$-complète,
% archive sheet number: 116
